\emph{Himpunan} yaiku kumpulan objèk kang diarani anggota himpunan.
Kita nulis $x \in A$ yèn $x$ dadi anggota~$A$. Himpunan iku
\emph{ekstensional}: ditemtokaké kanthi pepak déning anggotané.
Himpunan bisa diwènèhaké kanthi \emph{ndhaptar} anggotané siji-siji,
utawa kanthi nyebut sipat kang diduwèni bebarengan déning kabèh
anggotané (\emph{abstraksi}). Amarga ekstensionalitas, urutan lan
cara ndhaptar utawa nemtokaké anggota ora ngowahi himpunané.
Kanggo mbuktèkaké yèn $A$ lan~$B$ iku himpunan kang padha ($A=B$),
buktèkna yèn saben anggota~$A$ dadi anggota~$B$, lan kosok baliné.

Himpunan wigati kalebu wilangan asli ($\Nat$), wutuh ($\Int$),
rasional ($\Rat$), lan riil ($\Real$), uga \emph{rerangken}
($X^*$) lan \emph{urutan} tanpa wates ($X^\omega$) saka objèk.
$A$ iku \emph{subhimpunan} saka~$B$, $A\subseteq B$, yèn saben
anggota~$A$ uga dadi anggota~$B$. Kumpulan kabèh subhimpunan
saka~$B$ uga minangka himpunan, yaiku \emph{himpunan pangkat}
$\Pow{B}$. Kita bisa mbentuk \emph{gabungan} $A\cup B$ lan
\emph{irisan} $A\cap B$. \emph{Pasangan urut} $\tuple{x,y}$
ngemot rong objèk $x$ lan~$y$ kanthi urutan iku. Pasangan
$\tuple{x,y}$ lan $\tuple{y,x}$ béda kajaba yèn $x=y$.
Himpunan kabèh $\tuple{x,y}$ kanthi $x\in A$ lan $y\in B$
diarani \emph{prodhuk Kartesius} $A\times B$. Kita nulis $A^2$
kanggo $A\times A$. Upamané, $\Nat^2$ yaiku himpunan pasangan
wilangan asli.
